\chapter{射线与物质的相互作用}

\problem{8.1}{$^{137}\mathrm{Cs}$ 放射源发出的 $\gamma$ 射线能量是 $662\ \mathrm{keV}$，当它通过高纯锗半导体探测器时，试计算：
\begin{enumerate}
\item 光电子的能量；
\item 康普顿散射光子和反冲电子的最大和最小能量；
\item 连续两次 $90^\circ$ 康普顿散射光子的能量；
\item 激活的 K X 射线的能量。
\end{enumerate}}
% layout
\answerarea

\problem{8.2}{在一个用 X 射线激发的 X 荧光分析测量中，X 射线源为同位素 $^{238}\mathrm{Pu}$，放出的主要两条 X 射线的能量是 $17.3\ \mathrm{keV}$ 和 $13.6\ \mathrm{keV}$，用 Si(Li) 探测器测量样品发出的 X 射线。探测器和源相对样品来说放在同一边，可以近似地认为产生的 X 射线方向与源发出的 X 射线方向相反，Si(Li) 不能直接测到源发出的 X 射线。测量结果表明除了有样品的特征 X 射线外，还有 4 个峰，它们是怎样产生的？计算它们的能量。}
% layout
\answerarea
\newpage

\problem{8.3}{试求波长为 $337.1\ \mathrm{nm}$ 的氮分子激光束与静止电子以及和 $100\ \mathrm{keV}$ 能量的电子对头撞后的反散射光子的能量。}
% layout
\answerarea

\problem{8.4}{慢中子束通过一块镉板后强度减弱到原来的 $5\%$，镉的 $A=112$，$\rho=8.7\times10^3\ \mathrm{kg\,m^{-3}}$，中子是一个一个地被吸收，镉的吸收截面为 $2500\ \mathrm b$，求所需镉的厚度。}
% layout
\answerarea
\newpage

\problem{8.5}{已知铜和锌的 $\mathrm K_\alpha$ X 射线的波长分别为 $0.1539\ \mathrm{nm}$ 和 $0.1434\ \mathrm{nm}$，镍的 K 吸收限为 $0.1489\ \mathrm{nm}$，它对铜和锌的 $\mathrm K_\alpha$ 射线的质量吸收系数分别为 $0.480$ 和 $32.5\ \mathrm{m^2\,kg^{-1}}$，试问：
\begin{enumerate}
\item 为使黄铜的这两条 $\mathrm K_\alpha$ X 射线的相对强度之比提高 10 倍，需用多厚的镍吸收片？
\item 通过此吸收片后两条射线各减弱到多少？（已知镍的密度为 $9.04\times10^3\ \mathrm{kg\,m^{-3}}$。）
\end{enumerate}}
% layout
\answerarea

\problem{8.6}{在使用 X 光机和 $\gamma$ 放射源时，常需用铅板来屏蔽，以达到防护的目的，求为使 $^{137}\mathrm{Cs}$ 源的 $\gamma$ 射线（$662\ \mathrm{keV}$）和 $50\ \mathrm{keV}$ X 光机的 X 射线减弱到 $1/1000$ 所需的铅厚。（已知它们的质量吸收系数分别为 $0.010\ \mathrm{m^2\,kg^{-1}}$ 和 $0.849\ \mathrm{m^2\,kg^{-1}}$。）}
% layout
\answerarea
\newpage

\problem{8.7}{试证明正负电子对不能被单个孤立光子产生，只能发生在粒子附近。求光子在原子核或核外电子附近产生电子对效应的阈能。}
% layout
\answerarea

\problem{8.8}{泡室是一种径迹探测器，能把带电粒子的径迹拍成照片。一个光子通过泡室时产生了一个电子--正电子对，泡室放在强度 $B=0.20\ \mathrm T$ 的均匀磁场中，从照片中发现它们具有相反的曲率半径 $r=2.5\times10^{-2}\ \mathrm m$，求此光子的能量。}
% layout
\answerarea
\newpage

\problem{8.9}{在一束能量为 $10.0\ \mathrm{eV}$ 的电子射入单原子气体后，气体发生下列波长的光：$140.2\ \mathrm{nm}$、$253.6\ \mathrm{nm}$ 和 $313.2\ \mathrm{nm}$，波长为 $253.6\ \mathrm{nm}$ 的光强显著地大于另外两个波长的光强，试问出射电子具有多大能量？}
% layout
\answerarea

\problem{8.10}{试求为进行锌的 K X 荧光分析所需要的最小 X 射线、电子和质子的能量。}
% layout
\answerarea
\newpage

\problem{8.11}{试求能量为 $5.3\ \mathrm{MeV}$ 的 $\alpha$ 粒子
\begin{enumerate}
\item 穿过铝和铅的电离损失之比；
\item 穿过铝的电离损失值；
\item 和同样能量的氘核穿过铝的电离损失之比。
\end{enumerate}}
% layout
\answerarea

\problem{8.12}{在一个充 $1\ \mathrm{atm}$ 空气的云室内，$5.3\ \mathrm{MeV}$ 的 $\alpha$ 粒子显示出 $4\ \mathrm{cm}$ 长的径迹，问大约要观测多少个径迹才可能有一次机会发现由于核碰撞引起的大角弯折？}
% layout
\answerarea
\newpage

\problem{8.13}{求在液氮温度下处于三个最低转动态的 HD 分子的相对数。}
% layout
\answerarea

\problem{8.14}{求温度为 $3000\ \mathrm K$ 时处在基态、第一激发态和第二激发态的氢原子数目之比。}
% layout
\answerarea
\newpage

\problem{8.15}{原子在寿命为 $1.5\times10^{-8}\ \mathrm s$ 和 $2\times10^{-8}\ \mathrm s$ 的两激发态之间跃迁，产生相应于波长为 $532.0\ \mathrm{nm}$ 的谱线，试确定谱线的自然宽度。}
% layout
\answerarea

\problem{8.16}{在加热温度为 $440\ {}^\circ\mathrm C$ 时钠蒸气压为 $1\ \mathrm{mmHg}$。
\begin{enumerate}
\item 求此条件下钠蒸气原子共振谱线的自然宽度、碰撞展宽和多普勒展宽（已知钠的 $3\mathrm P$ 能级寿命为 $16\ \mathrm{ns}$）；
\item 能否观测到从 $4\mathrm D\to3\mathrm P$ 的 $568.822\ \mathrm{nm}$ 的发射和吸收谱线？
\end{enumerate}}
% layout
\answerarea
\newpage

\problem{8.17}{氢的 $2\mathrm P$ 能级的电偶极辐射的自发辐射寿命是 $1.596\ \mathrm{ns}$，试估计单电离氦和双电离锂的 $2\mathrm P$ 能级寿命。}
% layout
\answerarea

\newpage

